\documentclass[11pt,a4paper,spanish]{book}
\usepackage{estilo_unir-1}
\hypersetup{
	colorlinks   = true,
	urlcolor     = blue,
	linkcolor    = black,
	citecolor    = red
}
\usepackage[nottoc]{tocbibind}
\usepackage{listings}
\usepackage{graphicx}
\usepackage{longtable}
\usepackage{booktabs}
\lstset{
	basicstyle=\ttfamily\footnotesize,
	breaklines=true,
	frame=single,
	backgroundcolor=\color[gray]{0.95},
	numbers=left,
	numberstyle=\tiny\color[gray]{0.5},
	keywordstyle=\color{blue},
	commentstyle=\color[gray]{0.5},
	tabsize=2
}

%---------------------------
% Datos del trabajo
%---------------------------
\subject{Pr\'{a}cticas Acad\'{e}micas Externas}
\title{Sistema de medici\'{o}n de impedancia (EIS) para cuantificaci\'{o}n de Candida albicans --- Manual de usuario}
\author{Andr\'{e}s Alexandro Tapia Flores}
\profesor{Universidad Internacional de la Rioja (UNIR)}
\date{Mayo 2026}

\begin{document}
\renewcommand{\listfigurename}{\'{I}ndice de figuras}
\renewcommand{\listtablename}{\'{I}ndice de tablas}
\renewcommand{\contentsname}{\'{I}ndice de contenidos}
\renewcommand{\figurename}{Figura}
\renewcommand{\tablename}{Tabla}
\setlength{\headheight}{26pt}

\maketitle
\frontmatter
\tableofcontents
\listoffigures
\listoftables

\chapter{Introducci\'{o}n}

Este manual de usuario describe la instalaci\'{o}n, configuraci\'{o}n y operaci\'{o}n del sistema de medici\'{o}n de impedancia electroqu\'{i}mica (EIS) para la cuantificaci\'{o}n de \textit{Candida albicans}. El manual est\'{a} dirigido a operadores del sistema que necesiten realizar mediciones EIS, interpretar los resultados y exportar los datos para su an\'{a}lisis posterior.

Para los detalles t\'{e}cnicos de la arquitectura, el desarrollo del software y la implementaci\'{o}n del clasificador cu\'{a}ntico variacional, cons\'{u}ltense los entregables de las fases anteriores y la memoria final del proyecto.

\mainmatter

\chapter{Requisitos del sistema}

\section{Hardware}

Para operar con el equipo f\'{i}sico se requiere:
\begin{itemize}
	\item Microcontrolador ESP32 DevKit V1.
	\item Circuito integrado AD5933 con interfaz I\textsuperscript{2}C (o tarjeta de evaluaci\'{o}n EVAL-AD5940ELCZ).
	\item Celda electroqu\'{i}mica de 4 electrodos.
	\item Resistencia de calibraci\'{o}n de referencia (t\'{i}picamente 10~k$\Omega$).
	\item Cable USB para conexi\'{o}n serial entre ESP32 y PC.
\end{itemize}

Para operar en modo simulado (validaci\'{o}n sin hardware) solo se requiere el PC con el entorno software configurado.

\section{Software}

\begin{table}[ht]
\centering
\begin{tabular}{l l p{7cm}}
\hline\hline
Componente & Versi\'{o}n m\'{i}nima & Prop\'{o}sito \\ [0.5ex]
\hline
Python & 3.11 & Backend, collector, clasificador, coordinator \\
Node.js & 20 & Frontend React y herramientas de build \\
npm & 10 & Gestor de paquetes del frontend \\
uv & 0.11 & Gestor de dependencias Python \\
Navegador web & Chromium/Edge/Firefox & Visualizaci\'{o}n del dashboard \\
\hline
\end{tabular}
\caption[Requisitos software]{Requisitos software del sistema EIS. Fuente: elaboraci\'{o}n propia.}
\label{tab:requisitos-software}
\end{table}

\chapter{Instalaci\'{o}n}

\section{Clonar o descomprimir el proyecto}

El c\'{o}digo fuente se distribuye como un repositorio Git. Sit\'{u}ese en el directorio ra\'{i}z del proyecto:

\begin{lstlisting}[language=bash, caption={Estructura del proyecto}]
bkc32/
  src/              # Codigo fuente Python
  webapp/           # Aplicacion React
  documents/        # Entregables LaTeX
  data/             # Sesiones persistidas (generado en ejecucion)
  Makefile          # Comandos de automatizacion
  .env.example      # Variables de entorno de ejemplo
\end{lstlisting}

\section{Instalar dependencias Python}

El proyecto utiliza \texttt{uv} para la gesti\'{o}n de dependencias Python. Ejecute:

\begin{lstlisting}[language=bash]
uv sync
\end{lstlisting}

Esto instalar\'{a} \texttt{pyserial}, \texttt{websockets}, \texttt{python-dotenv} y dem\'{a}s dependencias listadas en \texttt{pyproject.toml}.

\section{Instalar dependencias del frontend}

\begin{lstlisting}[language=bash]
cd webapp
npm install
\end{lstlisting}

\section{Variables de entorno}

Copie el archivo de ejemplo y aj\'{u}stelo si es necesario:

\begin{lstlisting}[language=bash]
cp .env.example .env
\end{lstlisting}

En modo simulado no es necesario modificar \texttt{.env}, ya que el comando \texttt{make server-sim} utiliza el puerto serial virtual \texttt{/tmp/bkc32-sim-serial} por defecto.

\chapter{Inicio del sistema}

\section{Modo simulado (validaci\'{o}n sin hardware)}

Abra tres terminales y ejecute en cada una:

\textbf{Terminal 1} --- Tarjeta simulada:
\begin{lstlisting}[language=bash]
make sim-board
\end{lstlisting}

\textbf{Terminal 2} --- Backend:
\begin{lstlisting}[language=bash]
make server-sim
\end{lstlisting}

\textbf{Terminal 3} --- Frontend:
\begin{lstlisting}[language=bash]
make webapp
\end{lstlisting}

A continuaci\'{o}n abra el navegador en \url{http://localhost:5173}.

\section{Modo producci\'{o}n (con hardware ESP32/AD5933)}

Conecte el ESP32 al PC mediante USB. Identifique el puerto serial (por ejemplo, \texttt{/dev/ttyUSB0} en Linux o \texttt{COM3} en Windows) y config\'{u}relo en \texttt{.env}:

\begin{lstlisting}[language=bash]
SERIAL_PORT=/dev/ttyUSB0
BAUDRATE=115200
\end{lstlisting}

Luego inicie solo el backend y el frontend:

\textbf{Terminal 1} --- Backend:
\begin{lstlisting}[language=bash]
make server
\end{lstlisting}

\textbf{Terminal 2} --- Frontend:
\begin{lstlisting}[language=bash]
make webapp
\end{lstlisting}

\chapter{Flujo de operaci\'{o}n}

La Figura~\ref{fig:ui-estado-inicial} muestra el estado inicial de la interfaz tras establecer la conexi\'{o}n con el backend.

\begin{figure}[ht]
\centering
\includegraphics[width=\textwidth]{img/01_estado_inicial.png}
\caption[Estado inicial de la interfaz]{Estado inicial del dashboard con conexi\'{o}n activa, gr\'{a}ficas vac\'{i}as y workflow pendiente. Fuente: elaboraci\'{o}n propia.}
\label{fig:ui-estado-inicial}
\end{figure}

\section{Paso 1: Configurar el barrido}

En el panel \textit{Control de adquisici\'{o}n}, ajuste los par\'{a}metros:

\begin{table}[ht]
\centering
\begin{tabular}{l l p{8cm}}
\hline\hline
Par\'{a}metro & Valor por defecto & Descripci\'{o}n \\ [0.5ex]
\hline
\texttt{f\_min} & 1000~Hz & Frecuencia m\'{i}nima del barrido \\ 
\texttt{f\_max} & 100000~Hz & Frecuencia m\'{a}xima del barrido \\ 
\texttt{Puntos} & 50 & Cantidad de frecuencias equiespaciadas logar\'{i}tmicamente \\ 
\texttt{Settling} & 15 ciclos & Ciclos de asentamiento antes de medir \\ 
\texttt{R\_cal} & 10000~$\Omega$ & Resistencia de calibraci\'{o}n de referencia \\ 
Etiqueta esperada & Candida~(1) & Clase objetivo de la muestra: Control~(0) o Candida~(1) \\ 
\hline
\end{tabular}
\caption[Par\'{a}metros de configuraci\'{o}n]{Par\'{a}metros del panel de control. Fuente: elaboraci\'{o}n propia.}
\label{tab:parametros-configuracion}
\end{table}

Pulse \textbf{Configurar}. La interfaz mostrar\'{a} una notificaci\'{o}n de confirmaci\'{o}n como en la Figura~\ref{fig:ui-config-aplicada}.

\begin{figure}[ht]
\centering
\includegraphics[width=\textwidth]{img/02_configuracion_aplicada.png}
\caption[Configuraci\'{o}n aplicada]{Interfaz tras aplicar configuraci\'{o}n con los par\'{a}metros por defecto. Fuente: elaboraci\'{o}n propia.}
\label{fig:ui-config-aplicada}
\end{figure}

\section{Paso 2: Calibrar}

Pulse \textbf{Calibrar}. El sistema calcular\'{a} el factor de ganancia y el desfase de referencia utilizando la resistencia \texttt{R\_cal}. El panel mostrar\'{a} los valores obtenidos (Figura~\ref{fig:ui-calibracion-completada}).

\begin{figure}[ht]
\centering
\includegraphics[width=\textwidth]{img/03_calibracion_completada.png}
\caption[Calibraci\'{o}n completada]{Interfaz tras completar la calibraci\'{o}n. El workflow queda listo para iniciar el barrido. Fuente: elaboraci\'{o}n propia.}
\label{fig:ui-calibracion-completada}
\end{figure}

\textit{Nota:} Si modifica la configuraci\'{o}n despu\'{e}s de calibrar, la calibraci\'{o}n anterior quedar\'{a} invalidada. El sistema mostrar\'{a} una advertencia y ser\'{a} necesario recalibrar.

\section{Paso 3: Iniciar el barrido}

Seleccione la etiqueta esperada (Control o Candida) y pulse \textbf{Iniciar sweep}. El indicador de estado pasar\'{a} a \textsc{sweeping} y las gr\'{a}ficas se actualizar\'{a}n punto a punto conforme lleguen los datos por WebSocket (Figura~\ref{fig:ui-sweep-en-curso}).

\begin{figure}[ht]
\centering
\includegraphics[width=\textwidth]{img/04_sweep_en_curso.png}
\caption[Barrido en curso]{Interfaz durante la ejecuci\'{o}n de un barrido, con actualizaci\'{o}n en vivo de Bode y Nyquist. Fuente: elaboraci\'{o}n propia.}
\label{fig:ui-sweep-en-curso}
\end{figure}

\section{Paso 4: Revisar an\'{a}lisis}

Al finalizar el barrido, el panel de clasificaci\'{o}n mostrar\'{a}:
\begin{itemize}
	\item Probabilidad cu\'{a}ntica y etiqueta cu\'{a}ntica.
	\item Probabilidad cl\'{a}sica y etiqueta cl\'{a}sica.
	\item Coincidencia entre modelos y correspondencia con etiqueta esperada.
	\item Las seis caracter\'{i}sticas normalizadas utilizadas por ambos clasificadores.
\end{itemize}

La Figura~\ref{fig:ui-analisis} muestra el panel completo tras finalizar un barrido.

\begin{figure}[ht]
\centering
\includegraphics[width=\textwidth]{img/05_sweep_completado_analisis.png}
\caption[An\'{a}lisis al cierre del barrido]{Panel de clasificaci\'{o}n con probabilidades, features normalizadas, coincidencia entre modelos y rutas de exportaci\'{o}n. Fuente: elaboraci\'{o}n propia.}
\label{fig:ui-analisis}
\end{figure}

\section{Paso 5: Exportar datos}

Pulse \textbf{Exportar \'{u}ltimo} para generar los artefactos de la sesi\'{o}n m\'{a}s reciente. El panel \textit{Exportaci\'{o}n y metadatos} mostrar\'{a} las rutas de los archivos generados (Figura~\ref{fig:ui-exportacion}).

\begin{figure}[ht]
\centering
\includegraphics[width=\textwidth]{img/07_exportacion_generada.png}
\caption[Exportaci\'{o}n generada]{Panel de exportaci\'{o}n con rutas de CSV, metadata, summary y bundle. Fuente: elaboraci\'{o}n propia.}
\label{fig:ui-exportacion}
\end{figure}

Cada sesi\'{o}n genera autom\'{a}ticamente cuatro archivos en \texttt{data/acquisitions/<session\_id>/}:

\begin{table}[ht]
\centering
\begin{tabular}{l p{9cm}}
\hline\hline
Archivo & Contenido \\ [0.5ex]
\hline
\texttt{<id>\_data.csv} & Puntos EIS crudos con columnas \texttt{i,f,Z,phase,reZ,imZ,ts} \\
\texttt{<id>\_metadata.json} & Metadatos de sesi\'{o}n, eventos y resultado de clasificaci\'{o}n \\
\texttt{<id>\_summary.txt} & Resumen textual compacto para inspecci\'{o}n r\'{a}pida \\
\texttt{<id>\_bundle.json} & Paquete integral (metadatos + puntos + eventos) \\
\hline
\end{tabular}
\caption[Artefactos de exportaci\'{o}n]{Archivos generados por sesi\'{o}n. Fuente: elaboraci\'{o}n propia.}
\label{tab:artefactos-exportacion}
\end{table}

\section{Paso 6: Consultar historial}

Pulse \textbf{Historial} para cargar las sesiones previas almacenadas en disco. La tabla lateral muestra el identificador, n\'{u}mero de puntos, raz\'{o}n de cierre y probabilidades Q/C (Figura~\ref{fig:ui-historial}).

\begin{figure}[ht]
\centering
\includegraphics[width=\textwidth]{img/06_historial_sesiones.png}
\caption[Historial de sesiones]{Panel de historial con sesiones persistidas. Fuente: elaboraci\'{o}n propia.}
\label{fig:ui-historial}
\end{figure}

\chapter{Interpretaci\'{o}n de resultados}

\section{Diagramas de Bode}

\begin{itemize}
	\item \textbf{Bode magnitud} ($|Z|$ vs $f$, escala logar\'{i}tmica): la resistencia de la soluci\'{o}n $R_s$ se identifica como el plateau de alta frecuencia. La diferencia entre plateaus de baja y alta frecuencia corresponde a $R_{ct}$.
	\item \textbf{Bode fase} ($\phi$ vs $f$): $0^\circ$ indica comportamiento resistivo, $-90^\circ$ comportamiento puramente capacitivo. Los picos de fase corresponden a frecuencias donde dominan mecanismos de doble capa o transferencia de carga.
\end{itemize}

\section{Diagrama de Nyquist}

El semic\'{i}rculo corresponde al paralelo $R_{ct} \parallel C_{dl}$ con di\'{a}metro igual a $R_{ct}$. La intersecci\'{o}n izquierda con el eje real entrega $R_s$ y la derecha $R_s + R_{ct}$. La cola a $45^\circ$ a bajas frecuencias es la impedancia de Warburg asociada a difusi\'{o}n.

\section{Clasificaci\'{o}n}

El umbral de decisi\'{o}n es 0.5 para ambos clasificadores. La Tabla~\ref{tab:interpretacion-probabilidades} resume los rangos de interpretaci\'{o}n.

\begin{table}[ht]
\centering
\begin{tabular}{l c p{6.5cm}}
\hline\hline
Rango de probabilidad & Etiqueta & Interpretaci\'{o}n \\ [0.5ex]
\hline
0\%--20\% & 0 (Control) & Alta confianza de muestra control \\
20\%--40\% & 0 (Control) & Confianza moderada de control \\
40\%--50\% & 0 (Control) & Zona de incertidumbre, tendencia a control \\
50\%--60\% & 1 (Candida) & Zona de incertidumbre, tendencia a candida \\
60\%--80\% & 1 (Candida) & Confianza moderada de candida \\
80\%--100\% & 1 (Candida) & Alta confianza de candida \\
\hline
\end{tabular}
\caption[Interpretaci\'{o}n de probabilidades]{Interpretaci\'{o}n de los porcentajes mostrados por los clasificadores. Fuente: elaboraci\'{o}n propia.}
\label{tab:interpretacion-probabilidades}
\end{table}

\chapter{Resoluci\'{o}n de problemas}

\begin{table}[ht]
\centering
\begin{tabular}{p{4.5cm} p{9cm}}
\hline\hline
S\'{i}ntoma & Soluci\'{o}n \\ [0.5ex]
\hline
\textit{WebSocket desconectado} & Verifique que el backend est\'{e} ejecut\'{a}ndose y que no haya firewall bloqueando el puerto 8765. \\
\textit{No se puede iniciar sweep} & Aseg\'{u}rese de haber ejecutado \textbf{Calibrar} despu\'{e}s de \textbf{Configurar}. \\
\textit{Calibraci\'{o}n obsoleta} & Si cambia la configuraci\'{o}n despu\'{e}s de calibrar, recalibre antes de iniciar el sweep. \\
\textit{Sin datos en gr\'{a}ficas} & Verifique la conexi\'{o}n serial y que la tarjeta responda al comando PING. \\
\textit{Exportaci\'{o}n fallida} & Aseg\'{u}rese de que exista al menos una sesi\'{o}n finalizada antes de exportar. \\
\textit{Valores de impedancia err\'{a}ticos} & Verifique las conexiones de la celda electroqu\'{i}mica y la resistencia de calibraci\'{o}n. \\
\hline
\end{tabular}
\caption[Resoluci\'{o}n de problemas]{Gu\'{i}a de resoluci\'{o}n de problemas comunes. Fuente: elaboraci\'{o}n propia.}
\label{tab:resolucion-problemas}
\end{table}

\end{document}
